% GENERATED FILE. Edit canonical lessons, then run npm run generate:books.
% canonical-source-hash: fnv1a64:511633868aed9cb5
% canonical-lessons: PT-C276-quinzena, PT-C276-trimestre, PT-C276-mensal, PT-C276-milenio, PT-C276-posterior, PT-R276-revisao-269-276

\chapter{Quinzenas e trimestres}
\label{ch:pt-quinzenas-e-trimestres}
\hlchaptermodality{276}

\begin{chapteropening}
\textbf{By the end of this chapter:} \emph{I can count time in fortnights, terms, months and millennia, and say that something came later.}

Complete the last lesson of chapter 276: I can count time in fortnights, terms, months and millennia, and say that something came later.
\end{chapteropening}

\section[a quinzena]{a quinzena — the fortnight, two weeks}

\label{lesson:PT-C276-quinzena}



\begin{quote}
Before the new one: say the Portuguese for the twilight, then the Portuguese for the eve.
\end{quote}

\subsection*{You'll want to know: a quinzena}
\textbf{a quinzena} — \textquotedblleft{}the fortnight, two weeks\textquotedblright{}.

\begin{quote}
Passámos uma quinzena na praia.
\end{quote}

\begin{quote}
\emph{We spent a fortnight at the beach.}
\end{quote}

\emph{Na primeira quinzena de agosto} — in the first half of August.

\begin{cousinweb}
\emph{quinzena} is built on \emph{quinze} \textquotedblleft{}fifteen\textquotedblright{}, from Latin \emph{quīndecim} (Lewis \& Short): a stretch of fifteen days.
\end{cousinweb}

\begin{grammarlens}[title={What you've built}]
One more word for telling, explaining and describing at B1.
\end{grammarlens}

\subsection*{Your turn}
\begin{itemize}
\raggedright
  \item \emph{Say it:} \emph{a quinzena}
  \item \emph{Say it:} \emph{a quinzena}, once more
  \item \emph{Say it:} a sentence of your own with \emph{a quinzena}
  \item \emph{Recall:} say the Portuguese for the twilight, then the Portuguese for the eve, then say \emph{a quinzena} again
\end{itemize}

\begin{tcolorbox}[breakable,colback=teal!4,colframe=teal!35!black,title={Before you move on}]
What does a quinzena mean? (\textquotedblleft{}The fortnight, two weeks\textquotedblright{}.) Say it once more.
\end{tcolorbox}
\section[o trimestre]{o trimestre — the quarter, three months}

\label{lesson:PT-C276-trimestre}



\begin{quote}
Before the new one: say the Portuguese for the eve, then the Portuguese for the fortnight.
\end{quote}

\subsection*{You'll want to know: o trimestre}
\textbf{o trimestre} — \textquotedblleft{}the quarter, three months\textquotedblright{}.

\begin{quote}
A empresa vendeu muito no primeiro trimestre.
\end{quote}

\begin{quote}
\emph{The company sold a lot in the first quarter.}
\end{quote}

Firms count their year in \emph{trimestres}.

\begin{cousinweb}
\emph{trimestre} is a learned borrowing from Latin \emph{trimestris} \textquotedblleft{}of three months\textquotedblright{} (EtymDB), from \emph{trēs} \textquotedblleft{}three\textquotedblright{} and \emph{mēnsis} \textquotedblleft{}month\textquotedblright{}.
\end{cousinweb}

\begin{grammarlens}[title={What you've built}]
One more word for telling, explaining and describing at B1.
\end{grammarlens}

\subsection*{Your turn}
\begin{itemize}
\raggedright
  \item \emph{Say it:} \emph{o trimestre}
  \item \emph{Say it:} \emph{o trimestre}, once more
  \item \emph{Say it:} a sentence of your own with \emph{o trimestre}
  \item \emph{Recall:} say the Portuguese for the eve, then the Portuguese for the fortnight, then say \emph{o trimestre} again
\end{itemize}

\begin{tcolorbox}[breakable,colback=teal!4,colframe=teal!35!black,title={Before you move on}]
What does o trimestre mean? (\textquotedblleft{}The quarter, three months\textquotedblright{}.) Say it once more.
\end{tcolorbox}
\section[mensal]{mensal — monthly}

\label{lesson:PT-C276-mensal}



\begin{quote}
Before the new one: say the Portuguese for the fortnight, then the Portuguese for the quarter.
\end{quote}

\subsection*{You'll want to know: mensal}
\textbf{mensal} — \textquotedblleft{}monthly\textquotedblright{}.

\begin{quote}
Pago uma renda mensal de quinhentos euros.
\end{quote}

\begin{quote}
\emph{I pay a monthly rent of five hundred euros.}
\end{quote}

Built on \emph{mês}. \emph{Semanal} — weekly; \emph{anual} — yearly.

\begin{cousinweb}
\emph{mensal} is a learned word built on Latin \emph{mēnsis} \textquotedblleft{}a month\textquotedblright{} (Lewis \& Short), with the adjective ending \emph{-al}.
\end{cousinweb}

\begin{grammarlens}[title={What you've built}]
One more word for telling, explaining and describing at B1.
\end{grammarlens}

\subsection*{Your turn}
\begin{itemize}
\raggedright
  \item \emph{Say it:} \emph{mensal}
  \item \emph{Say it:} \emph{mensal}, once more
  \item \emph{Say it:} a sentence of your own with \emph{mensal}
  \item \emph{Recall:} say the Portuguese for the fortnight, then the Portuguese for the quarter, then say \emph{mensal} again
\end{itemize}

\begin{tcolorbox}[breakable,colback=teal!4,colframe=teal!35!black,title={Before you move on}]
What does mensal mean? (\textquotedblleft{}Monthly\textquotedblright{}.) Say it once more.
\end{tcolorbox}
\section[o milénio]{o milénio — the millennium}

\label{lesson:PT-C276-milenio}



\begin{quote}
Before the new one: say the Portuguese for the quarter, then the Portuguese for monthly.
\end{quote}

\subsection*{You'll want to know: o milénio}
\textbf{o milénio} — \textquotedblleft{}the millennium\textquotedblright{}.

\begin{quote}
Esta igreja tem quase um milénio.
\end{quote}

\begin{quote}
\emph{This church is almost a thousand years old.}
\end{quote}

In Brazil: \emph{milênio}.

\begin{cousinweb}
\emph{milénio} is a learned borrowing from Latin \emph{millennium} (EtymDB), a word made from \emph{mīlle} \textquotedblleft{}a thousand\textquotedblright{} and \emph{annus} \textquotedblleft{}year\textquotedblright{}.
\end{cousinweb}

\begin{grammarlens}[title={What you've built}]
One more word for telling, explaining and describing at B1.
\end{grammarlens}

\subsection*{Your turn}
\begin{itemize}
\raggedright
  \item \emph{Say it:} \emph{o milénio}
  \item \emph{Say it:} \emph{o milénio}, once more
  \item \emph{Say it:} a sentence of your own with \emph{o milénio}
  \item \emph{Recall:} say the Portuguese for the quarter, then the Portuguese for monthly, then say \emph{o milénio} again
\end{itemize}

\begin{tcolorbox}[breakable,colback=teal!4,colframe=teal!35!black,title={Before you move on}]
What does o milénio mean? (\textquotedblleft{}The millennium\textquotedblright{}.) Say it once more.
\end{tcolorbox}
\section[posterior]{posterior — subsequent, later; back, rear}

\label{lesson:PT-C276-posterior}



\begin{quote}
Before the new one: say the Portuguese for monthly, then the Portuguese for the millennium.
\end{quote}

\subsection*{You'll want to know: posterior}
\textbf{posterior} — \textquotedblleft{}subsequent, later; back, rear\textquotedblright{}.

\begin{quote}
O médico marcou uma consulta posterior.
\end{quote}

\begin{quote}
\emph{The doctor booked a later appointment.}
\end{quote}

Also ‘at the back’: \emph{a parte posterior do carro}.

\begin{cousinweb}
\emph{posterior} is a learned borrowing from Latin \emph{posterior} \textquotedblleft{}coming after, later\textquotedblright{} (EtymDB), the comparative of \emph{posterus}, from \emph{post} \textquotedblleft{}after\textquotedblright{}.
\end{cousinweb}

\begin{grammarlens}[title={What you've built}]
One more word for telling, explaining and describing at B1.
\end{grammarlens}

\subsection*{Your turn}
\begin{itemize}
\raggedright
  \item \emph{Say it:} \emph{posterior}
  \item \emph{Say it:} \emph{posterior}, once more
  \item \emph{Say it:} a sentence of your own with \emph{posterior}
  \item \emph{Recall:} say the Portuguese for monthly, then the Portuguese for the millennium, then say \emph{posterior} again
\end{itemize}

\begin{tcolorbox}[breakable,colback=teal!4,colframe=teal!35!black,title={Before you move on}]
What does posterior mean? (\textquotedblleft{}Subsequent, later; back, rear\textquotedblright{}.) Say it once more.
\end{tcolorbox}
\section[(dialogue)]{Revisão: quando e como}

\label{lesson:PT-R276-revisao-269-276}



\begin{quote}
Say the Portuguese for firstly and initially.
\end{quote}

\subsection*{Your turn}
Say these aloud:

\begin{itemize}
\raggedright
  \item firstly, initially, next, later, finally
  \item meanwhile, simultaneously, as soon as, the moment that, suddenly
  \item to occur, to come after, to take place, the event, the incident
  \item to last, to prolong, the instant, the period, immediately
  \item the stage, the phase, the sequence, the succession, the chronology
  \item formerly, lately, after all, gradually, little by little
  \item the dawn, the late afternoon, nightfall, the twilight, the eve
  \item the fortnight, the quarter, monthly, the millennium, subsequent
\end{itemize}

\begin{tcolorbox}[breakable,colback=teal!4,colframe=teal!35!black,title={Before you move on}]
Firstly? (\textbf{primeiramente}.) Initially? (\textbf{inicialmente}.) Next? (\textbf{seguidamente}.) Later? (\textbf{posteriormente}.) Finally? (\textbf{por fim}.) Meanwhile? (\textbf{entretanto}.) Simultaneously? (\textbf{simultaneamente}.) As soon as? (\textbf{assim que}.) The moment that? (\textbf{logo que}.) Suddenly? (\textbf{subitamente}.) To occur? (\textbf{ocorrer}.) To come after? (\textbf{suceder}.) To take place? (\textbf{decorrer}.) The event? (\textbf{o acontecimento}.) The incident? (\textbf{a ocorrência}.) To last? (\textbf{durar}.) To prolong? (\textbf{prolongar}.) The instant? (\textbf{o instante}.) The period? (\textbf{o período}.) Immediately? (\textbf{imediatamente}.) The stage? (\textbf{a etapa}.) The phase? (\textbf{a fase}.) The sequence? (\textbf{a sequência}.) The succession? (\textbf{a sucessão}.) The chronology? (\textbf{a cronologia}.) Formerly? (\textbf{outrora}.) Lately? (\textbf{ultimamente}.) After all? (\textbf{afinal}.) Gradually? (\textbf{gradualmente}.) Little by little? (\textbf{aos poucos}.) The dawn? (\textbf{o amanhecer}.) The late afternoon? (\textbf{o entardecer}.) Nightfall? (\textbf{o anoitecer}.) The twilight? (\textbf{o crepúsculo}.) The eve? (\textbf{a véspera}.) The fortnight? (\textbf{a quinzena}.) The quarter? (\textbf{o trimestre}.) Monthly? (\textbf{mensal}.) The millennium? (\textbf{o milénio}.) Subsequent? (\textbf{posterior}.)
\end{tcolorbox}
